% =============================================================================
%  Joint Neural Architecture and Quantization Search for Real-Time Anomaly
%  Detection on Embedded Manipulators
%
%  LaTeX draft — Sections writable from current experimental data.
%  Placeholders marked with \TODO{} require Jetson deployment results.
%  Intended format: IEEE Transactions on Industrial Informatics (double-column).
% =============================================================================

\documentclass[journal]{IEEEtran}

% ---- Packages ---------------------------------------------------------------
\usepackage{amsmath,amssymb,amsfonts}
\usepackage{algorithmic}
\usepackage{algorithm}
\usepackage{array}
\usepackage{booktabs}
\usepackage{multirow}
\usepackage{graphicx}
\usepackage{cite}
\usepackage{xcolor}
\usepackage{hyperref}
\usepackage{bm}
\usepackage{bbm}

% ---- Convenience macros -----------------------------------------------------
\newcommand{\TODO}[1]{\textcolor{red}{\textbf{[TODO: #1]}}}
\newcommand{\Rbb}{\mathbb{R}}
\newcommand{\Nbb}{\mathbb{N}}
\newcommand{\calA}{\mathcal{A}}
\newcommand{\calQ}{\mathcal{Q}}
\newcommand{\calX}{\mathcal{X}}
\newcommand{\calF}{\mathcal{F}}
\newcommand{\argmin}{\operatorname*{arg\,min}}

% =============================================================================
\begin{document}

\title{Joint Neural Architecture and Quantization Search\\
for Real-Time Anomaly Detection on Embedded Manipulators}

% Author block — fill in before submission
\author{Maicol~Cifuentes,~\IEEEmembership{Student Member,~IEEE}%
\thanks{M. Cifuentes is with the Universidad Tecnol\'ogica de Pereira (UTP),
Pereira, Colombia (e-mail: mcifuentes@utp.edu.co).}%
\thanks{Manuscript received \TODO{date}; revised \TODO{date}.}}

\maketitle

% =============================================================================
\begin{abstract}
Deploying vision-based anomaly detection inside robotic manipulators demands
models that are simultaneously accurate, low-latency, memory-efficient, and
energy-frugal — objectives that conflict in all known neural architectures.
We present a joint Neural Architecture Search and Quantization-Aware Training
(NAS+QAT) framework that resolves this conflict through a four-objective
Pareto optimization.
A genome representation encodes both the macro-structure of the detector
(family, depth, channel widths, block types, attention mechanisms) and its
per-layer mixed-precision quantization profile (weight and activation bit
widths independently drawn from $\{4,6,8\}$ bits), yielding a joint search
space with more than $10^{10}$ candidate configurations.
NSGA-II is used to evolve a population of 32 architectures over 30 generations,
minimizing simultaneously the complement of the image-level Area Under the
ROC Curve (AUROC), inference latency, peak RAM, and energy per inference.
Experiments on the MVTec Anomaly Detection benchmark demonstrate that the
evolutionary search identifies 32 non-dominated solutions spanning AUROC from
0.781 to 0.894 at proxy evaluation, with proxy latencies between 12.0 and
14.7\,ms.
Full-budget Quantization-Aware Training of the five highest-ranked Pareto
candidates consistently raises AUROC into the range 0.900--0.914, yielding an
average accuracy gain of $+2.4$ percentage points over their proxy estimates.
Notably, one candidate — a depth-5 UNet with selective CBAM attention and
aggressive mixed-precision quantization — reaches an AUROC of 0.914 while
converging to its best checkpoint three times faster than comparable models
(epoch 48 vs.\ 185 on average).
\TODO{After TensorRT compilation and benchmarking on a Jetson Orin Nano 8\,GB,
the top-ranked model achieves XX\,ms latency (p95), XX\,mJ per inference, and
XX\,MB peak RAM in INT8 precision, with an on-device AUROC drop below the
$0.02$ tolerance threshold.}
The resulting framework provides a principled, reproducible methodology for
designing task-specific quantized detectors targeting embedded robotic
platforms.
\end{abstract}

\begin{IEEEkeywords}
anomaly detection, neural architecture search, quantization-aware training,
mixed-precision, multi-objective optimization, NSGA-II, embedded AI,
robotic manipulator, MVTec AD, Jetson Orin Nano.
\end{IEEEkeywords}

% =============================================================================
\section{Introduction}
\label{sec:intro}

Industrial robotic manipulators are increasingly expected to perform in-line
quality inspection alongside their primary manipulation tasks.
Embedding an anomaly detector directly on the robot's on-board compute
eliminates the latency and bandwidth overhead of streaming images to a remote
server, enables closed-loop rejection of defective parts, and reduces
infrastructure cost.
The Jetson Orin Nano 8\,GB~\cite{nvidia2023jetson}, with its Ampere GPU
and unified memory architecture, is a representative platform for this
application: it provides sufficient compute for deep inference while fitting
within the power and size envelope of an articulated arm.

However, deploying state-of-the-art anomaly detectors on such hardware is
far from straightforward.
Methods that achieve near-perfect performance on the MVTec Anomaly Detection
(MVTec AD) benchmark~\cite{bergmann2019mvtec} — such as
PatchCore~\cite{roth2022towards}, CFlow-AD~\cite{gudovskiy2022cflow}, and
FastFlow~\cite{yu2021fastflow} — were designed for server-class GPUs and
typically require hundreds of megabytes of memory, dozens of milliseconds of
latency per frame, and watt-level power draws that are incompatible with
embedded operation.
Lightweight alternatives such as
EfficientAD~\cite{batzner2024efficientad} improve efficiency but still require
manual architecture selection and do not expose a systematic mechanism for
adapting the accuracy--efficiency trade-off to a given hardware budget.

A natural approach is to automate this adaptation via Neural Architecture Search
(NAS)~\cite{elsken2019neural,white2023neural}.
Yet existing NAS frameworks for anomaly detection treat architecture and
quantization as separate design decisions: architecture search selects a
floating-point model, and post-training quantization (PTQ) or
quantization-aware training (QAT) is applied afterward, often degrading
accuracy for aggressive bit widths.
We argue that architecture and quantization must be co-optimized: the optimal
per-layer bit assignment depends on the sensitivity of each layer's
representations, which in turn depends on architectural choices such as skip
connections, attention modules, and block type.

This paper addresses the following research question:
\textit{Can we jointly search the architecture and quantization configuration
of an anomaly detection network so that the resulting model simultaneously
satisfies accuracy, latency, memory, and energy requirements of an embedded
robotic platform, without requiring any target-device execution during the
search itself?}

Our answer is a \textbf{Joint NAS+QAT} framework structured around three
contributions:

\begin{enumerate}

\item \textbf{A joint search space} that unifies four anomaly-detection
architecture families (patch-based CNN, encoder--decoder UNet, student--teacher
distillation, and autoencoder) with a per-layer mixed-precision quantization
profile.
Each genome encodes $d$ architectural and $2d$ quantization decisions
independently for a network of depth $d$, yielding a search space of
$|\calX| > 10^{10}$ configurations.

\item \textbf{A four-objective NSGA-II search} that simultaneously optimizes
accuracy (AUROC), inference latency, peak RAM, and energy per inference using
only inexpensive proxy evaluations on a development PC, avoiding the need for
target-device access during the evolutionary search.
After 30 generations with a population of 32, the algorithm identifies 32
Pareto-optimal configurations spanning AUROC 0.781--0.894 at proxy-latency
12.0--14.7\,ms.

\item \textbf{Full-budget QAT retraining analysis} showing that the Pareto
ordering established by proxy metrics partially but not perfectly predicts
post-retraining accuracy: all five retrained candidates exceed AUROC\,0.90,
but re-ranking occurs, with the search's fourth-ranked candidate overtaking the
first after 200 epochs.
This re-ranking is accompanied by a convergence anomaly — one architecture
reaches its global minimum in 48 epochs versus an average of 185 for its
peers — which we attribute to its lower-layer bottleneck ratio and
selective use of CBAM attention.

\end{enumerate}

The remainder of the paper is organized as follows.
Section~\ref{sec:related} reviews related work on anomaly detection,
NAS, and network quantization.
Section~\ref{sec:formulation} formalizes the multi-objective optimization
problem.
Section~\ref{sec:searchspace} describes the joint search space and genome
encoding.
Section~\ref{sec:nsga2} details the NSGA-II search procedure and proxy
fitness evaluation.
Section~\ref{sec:retraining} describes the full-budget QAT retraining stage.
Section~\ref{sec:experiments} presents the experimental setup and results.
Section~\ref{sec:discussion} discusses limitations and open questions.
Section~\ref{sec:conclusion} concludes.

% =============================================================================
\section{Related Work}
\label{sec:related}

\subsection{Visual Anomaly Detection}

Unsupervised anomaly detection from images has been studied extensively in the
context of industrial quality control, driven by the MVTec AD
benchmark~\cite{bergmann2019mvtec} and its successors
MVTec LOCO~\cite{bergmann2022beyond} and
VisA~\cite{zou2022spot}.
Methods can be broadly categorized as reconstruction-based, embedding-based,
normalizing-flow-based, and knowledge-distillation-based.

\textit{Reconstruction-based} methods~\cite{bergmann2019improving,
baur2021autoencoders} train autoencoders on normal data and flag regions
where reconstruction error exceeds a threshold.
While simple to implement, they suffer from over-generalization: powerful
decoders can reconstruct anomalies accurately, suppressing the detection
signal.

\textit{Embedding-based} methods project image patches into a
pre-trained feature space and flag outliers by distance or density estimation.
SPADE~\cite{cohen2020sub} uses $k$-NN in feature space;
PaDiM~\cite{defard2021padim} models per-position Gaussian distributions;
PatchCore~\cite{roth2022towards} maintains a coreset of normal patch
embeddings and achieves near-human performance on MVTec AD.
These methods are accurate but memory-intensive — PatchCore requires storing
the entire coreset at inference time, making it unsuitable for embedded
deployment.

\textit{Normalizing-flow-based} methods~\cite{gudovskiy2022cflow,
yu2021fastflow,rudolph2022csflow} model the distribution of normal features
as an exact likelihood under an invertible transformation.
They avoid memory overhead but require significant compute at inference due
to the flow's forward pass.

\textit{Knowledge-distillation-based} methods~\cite{bergmann2020uninformed,
wang2021student,batzner2024efficientad} train a student network to mimic a
teacher on normal data; anomalies cause elevated discrepancy.
EfficientAD~\cite{batzner2024efficientad} is currently the most efficient
single-model detector, reaching sub-millisecond latency on a server GPU while
maintaining competitive AUROC, but it is not designed for integer quantization
or embedded execution.

To our knowledge, no prior work jointly optimizes the architecture and
quantization of an anomaly detector for a specific embedded platform.
The closest works are
EfficientAD~\cite{batzner2024efficientad}, which targets efficiency but not
quantization, and
NASiam~\cite{TODO_nasiam}, which applies NAS to anomaly detection but ignores
quantization and hardware constraints.

\subsection{Neural Architecture Search}

NAS automates the design of neural network topologies for a target task.
Early approaches based on reinforcement learning~\cite{zoph2017neural} and
evolutionary algorithms~\cite{real2019regularized} required thousands of GPU
hours.
Differentiable NAS (DARTS~\cite{liu2019darts} and its variants) reduced
search cost dramatically by relaxing the discrete architecture selection to a
continuous mixture of operations, but is prone to performance collapse and
produces architectures that do not transfer reliably to quantized or
resource-constrained settings.

\textit{Hardware-aware NAS} incorporates device-specific metrics into the
search objective.
MNASNet~\cite{tan2019mnasnet} uses a reward that multiplies accuracy by a
latency term; ProxylessNAS~\cite{cai2019proxylessnas} allows direct latency
measurement on target hardware during search;
Once-for-All~\cite{cai2020ofa} trains a super-network from which sub-networks
of varying sizes can be extracted without re-training.
These methods target classification tasks and assume full-precision inference.

\textit{Multi-objective NAS} explicitly maintains a Pareto front over
competing objectives.
NSGA-NET~\cite{lu2019nsga} applies NSGA-II~\cite{deb2002fast} to the
architecture search; CARS~\cite{yang2020cars} uses a co-evolutionary strategy;
LEMONADE~\cite{elsken2019lemonade} performs Lamarckian evolution with mutation
operators designed to preserve architectural patterns.
None of these methods include quantization in the search variable.

\subsection{Network Quantization and Mixed Precision}

Quantization reduces model size and accelerates inference by representing
weights and activations in low-bit integer formats.
Post-Training Quantization (PTQ) requires only a small calibration
dataset~\cite{nagel2020up,hubara2021accurate} but degrades accuracy
significantly below 8 bits.
QAT~\cite{jacob2018quantization} incorporates fake-quantization nodes during
training and uses the Straight-Through Estimator (STE)~\cite{bengio2013estimating}
to propagate gradients through the discontinuous round operation, typically
achieving accuracy within 1\% of the full-precision baseline even at 4 bits.

\textit{Mixed-precision} QAT assigns different bit widths to different layers,
exploiting the observation that some layers are far more sensitive to
quantization than others~\cite{dong2019hawq,wang2019haq}.
HAQ~\cite{wang2019haq} uses reinforcement learning to select per-layer bit
widths; HAWQ~\cite{dong2019hawq} uses the Hessian spectrum as a sensitivity
proxy.
APQ~\cite{wang2020apq} jointly searches architecture and quantization policy
using a once-for-all super-network, but restricts its evaluation to ImageNet
classification.
Our work differs in two key aspects: (i) we target anomaly detection, where
the task structure (one-class learning from normal data) fundamentally differs
from supervised classification; and (ii) we perform a \emph{from-scratch}
joint search without pre-training a super-network, which avoids the coupling
between the super-network's training and the final model's quantization
sensitivity.

\subsection{Embedded Deployment of Anomaly Detectors}

Deploying anomaly detectors on edge devices has received growing attention
as robotic systems become more autonomous.
Schlüter et al.~\cite{schluter2022natural} evaluate several MVTec AD methods
on a Raspberry Pi but restrict to unquantized float32 models.
Batzner et al.~\cite{batzner2024efficientad} benchmark EfficientAD on a
Jetson AGX Orin at full precision.
Quantized deployment of anomaly detectors — particularly with layer-wise
mixed precision targeting INT8 TensorRT engines — has not been systematically
studied, to the best of our knowledge.
This paper fills that gap.

% =============================================================================
\section{Problem Formulation}
\label{sec:formulation}

Let $\bm{x} = (\bm{\alpha}, \bm{q}) \in \calX = \calA \times \calQ$ denote a
candidate solution, where $\bm{\alpha} \in \calA$ is the architecture
specification and $\bm{q} \in \calQ$ is the quantization specification.
Given a training dataset $\mathcal{D}_{\mathrm{train}}$ of normal images and
a test set $\mathcal{D}_{\mathrm{test}}$ with ground-truth anomaly labels, we
define four objective functions:

\begin{align}
f_1(\bm{x}) &= 1 - \mathrm{AUROC}\bigl(\bm{x}, \mathcal{D}_{\mathrm{test}}\bigr)
  \label{eq:obj1} \\
f_2(\bm{x}) &= L(\bm{x})   \quad \text{[ms]}  \label{eq:obj2} \\
f_3(\bm{x}) &= M(\bm{x})   \quad \text{[MB]}  \label{eq:obj3} \\
f_4(\bm{x}) &= E(\bm{x})   \quad \text{[mJ]}  \label{eq:obj4}
\end{align}

where $\mathrm{AUROC}(\bm{x}, \cdot)$ is the image-level area under the
receiver operating characteristic curve, $L(\bm{x})$ is the per-image
inference latency at the 95th percentile, $M(\bm{x})$ is the peak GPU RAM
consumed during a single forward pass, and $E(\bm{x})$ is the energy consumed
per inference.
We seek to solve the four-objective minimization problem:

\begin{equation}
\min_{\bm{x} \in \calX}
\bigl[ f_1(\bm{x}),\; f_2(\bm{x}),\; f_3(\bm{x}),\; f_4(\bm{x}) \bigr]
\label{eq:mop}
\end{equation}

subject to the hardware budget constraints:

\begin{align}
f_2(\bm{x}) &\leq L_{\max} = 20\;\text{ms}  \label{eq:c1} \\
f_3(\bm{x}) &\leq M_{\max} = 512\;\text{MB} \label{eq:c2} \\
f_4(\bm{x}) &\leq E_{\max} = 5\;\text{mJ}   \label{eq:c3}
\end{align}

\noindent
Because there is in general no single solution that minimizes all four
objectives simultaneously, we target the \emph{Pareto front}: the set
$\mathcal{P}^* \subset \calX$ of all non-dominated solutions.
A solution $\bm{x}_1$ \textit{Pareto-dominates} $\bm{x}_2$, written
$\bm{x}_1 \prec \bm{x}_2$, if and only if

\begin{equation}
\forall i \in \{1,2,3,4\}: f_i(\bm{x}_1) \leq f_i(\bm{x}_2)
\quad\text{and}\quad
\exists j: f_j(\bm{x}_1) < f_j(\bm{x}_2).
\label{eq:dominance}
\end{equation}

\noindent
The Pareto front is:
$\mathcal{P}^* = \{\bm{x} \in \calX \mid \nexists\, \bm{x}' \in \calX : \bm{x}' \prec \bm{x}\}$.

\medskip
\noindent\textbf{Proxy vs.\ ground-truth objectives.}
Evaluating $f_2$, $f_3$, and $f_4$ exactly requires running inference on the
target Jetson hardware, which is expensive and limits the number of candidates
that can be evaluated during search.
We therefore introduce proxy functions
$\tilde{f}_2(\bm{x})$, $\tilde{f}_3(\bm{x})$, $\tilde{f}_4(\bm{x})$
that approximate the true objectives using only PC-side computation
(see Section~\ref{sec:proxy}).
The quality of this approximation is validated post-hoc in
Section~\ref{sec:deploy_results} \TODO{after on-device benchmarking}.

% =============================================================================
\section{Joint Search Space}
\label{sec:searchspace}

The joint search space $\calX = \calA \times \calQ$ decomposes into an
architecture component and a quantization component, described below.

\subsection{Architecture Search Space}
\label{sec:arch_space}

An architecture specification $\bm{\alpha}$ consists of:

\begin{itemize}
\item \textbf{Family} $\phi \in \{\texttt{patch\_cnn},\, \texttt{unet},\,
  \texttt{student\_teacher},\, \texttt{autoencoder}\}$: the high-level
  inductive bias of the detector.
  \texttt{patch\_cnn} computes patch-level anomaly scores directly;
  \texttt{unet} uses an encoder--decoder with skip connections for
  pixel-level reconstruction; \texttt{student\_teacher} distills features
  from a frozen teacher into a trained student;
  \texttt{autoencoder} learns a compact latent representation of normal
  appearance.
\item \textbf{Input resolution} $s \in \{224, 256\}$ pixels (square).
\item \textbf{Depth} $d \in \{2, 3, 4, 5\}$: number of encoder stages.
\item Per-stage parameters (vectors of length $d$):
  \begin{itemize}
    \item \textbf{Channel widths} $\bm{c} \in \{16, 24, 32, 48, 64, 96, 128, 192\}^d$.
    \item \textbf{Kernel sizes} $\bm{k} \in \{3, 5, 7\}^d$.
    \item \textbf{Block types} $\bm{b} \in \{\texttt{ir},\,
      \texttt{conv},\, \texttt{ds}\}^d$:
      inverted-residual (MobileNetV2-style), standard convolution, and
      depthwise-separable convolution, respectively.
    \item \textbf{Bottleneck ratios} $\bm{r} \in \{0.5, 1.0, 2.0, 4.0, 6.0\}^d$.
    \item \textbf{Skip connections} $\bm{s} \in \{0, 1\}^d$.
    \item \textbf{Attention mechanisms}
      $\bm{a} \in \{\texttt{none},\, \texttt{cbam},\, \texttt{eca},\,
      \texttt{se},\, \texttt{spatial}\}^d$.
  \end{itemize}
\end{itemize}

\noindent
The Convolutional Block Attention Module
(CBAM)~\cite{woo2018cbam} applies sequential channel-wise and spatial-wise
attention; the Efficient Channel Attention (ECA)~\cite{wang2020eca} uses a
1-D convolution over channel statistics; the Squeeze-and-Excitation
(SE)~\cite{hu2018squeeze} module applies channel-wise re-weighting; and
spatial attention applies a learned $7\times7$ spatial gate.

\subsection{Quantization Search Space}
\label{sec:quant_space}

A quantization specification $\bm{q}$ consists of:

\begin{itemize}
\item Per-layer weight bit widths $\bm{w} \in \{4, 6, 8\}^d$.
\item Per-layer activation bit widths $\bm{a} \in \{4, 6, 8\}^d$.
\item Symmetric quantization flag $s_q \in \{0, 1\}$:
  when enabled, the quantization range is $[-\Delta, \Delta]$ (signed
  symmetric); otherwise asymmetric unsigned.
\item Per-channel quantization flag $p_c \in \{0, 1\}$:
  when enabled, each output channel of a weight tensor uses an independent
  scale factor.
\end{itemize}

\noindent
The $n$-bit uniform affine quantization of a real-valued tensor $\bm{v}$ is:

\begin{equation}
Q_n(\bm{v}) = \mathrm{clip}\!\left(
  \left\lfloor \frac{\bm{v}}{s} \right\rceil + z,\;
  q_{\min},\; q_{\max}
\right),
\label{eq:quant}
\end{equation}

where $s > 0$ is the scale factor, $z \in \Nbb$ is the zero-point
(zero for symmetric), $q_{\min} = -(2^{n-1})$,
$q_{\max} = 2^{n-1} - 1$ for signed, and
$\lfloor \cdot \rceil$ denotes rounding to the nearest integer.
During QAT, the forward pass uses $Q_n$ while gradients are propagated
with the Straight-Through Estimator~\cite{bengio2013estimating}:

\begin{equation}
\frac{\partial \mathcal{L}}{\partial \bm{v}}
= \frac{\partial \mathcal{L}}{\partial Q_n(\bm{v})}
\cdot \mathbbm{1}\!\left[q_{\min} \leq \frac{\bm{v}}{s} + z \leq q_{\max}\right].
\label{eq:ste}
\end{equation}

\subsection{Genome Encoding}
\label{sec:genome}

Each candidate $\bm{x} = (\bm{\alpha}, \bm{q})$ is encoded as a fixed-length
integer vector $\bm{g} \in \mathcal{G}$, the genome.
For a network of depth $d$, the genome has dimension:

\begin{equation}
|\bm{g}| = \underbrace{4}_{\phi,\, s,\, d,\, p_c} +
\underbrace{6d}_{c_i, k_i, b_i, r_i, s_i, a_i} +
\underbrace{2d + 2}_{w_i, a_i, s_q, p_c},
\label{eq:genome_dim}
\end{equation}

which ranges from 22 (at $d=2$) to 46 (at $d=5$).
The genome is padded to length 46 for crossover and mutation operations that
require fixed-length representations.
Inactive positions (corresponding to stages beyond depth $d$) are masked
and do not affect phenotype decoding.

\noindent\textbf{Search space cardinality.}
Taking the Cartesian product of all discrete choices and integrating over all
depths, the total search space contains:

\begin{equation}
|\calX| \;=\; \sum_{d=2}^{5}
  |\calA_d| \cdot |\calQ_d|
  \;\approx\; 2.3 \times 10^{10}
\label{eq:space_size}
\end{equation}

configurations, making exhaustive evaluation intractable and motivating the
use of evolutionary search.

% =============================================================================
\section{Multi-Objective Evolutionary Search}
\label{sec:nsga2}

\subsection{NSGA-II Overview}

We adopt the Non-dominated Sorting Genetic Algorithm II
(NSGA-II)~\cite{deb2002fast} as the outer optimizer.
NSGA-II maintains a population $\mathcal{P}_t$ of $N$ individuals at each
generation $t$, applies selection, crossover, and mutation to produce $N$
offspring, and combines parents and offspring into a pool of size $2N$ from
which the next generation is selected by non-dominated sorting and crowding
distance.
The overall pipeline is given in Algorithm~\ref{alg:nsga2}.

\begin{algorithm}
\caption{Joint NAS+QAT Search via NSGA-II}
\label{alg:nsga2}
\begin{algorithmic}[1]
\REQUIRE Population size $N$, generations $T$,
  training set $\mathcal{D}_{\mathrm{train}}$, proxy evaluator $\tilde{\mathcal{F}}$.
\ENSURE Approximate Pareto front $\hat{\mathcal{P}}$.
\STATE \textbf{Initialize}: $\mathcal{P}_0 \leftarrow$ random($N$) \COMMENT{uniform over $\calX$}
\FOR{$t = 0$ \TO $T-1$}
  \FOR{each $\bm{x} \in \mathcal{P}_t$ not yet evaluated}
    \STATE $\bm{f}(\bm{x}) \leftarrow \tilde{\mathcal{F}}(\bm{x}, \mathcal{D}_{\mathrm{train}})$
    \COMMENT{proxy fitness; see \S\ref{sec:proxy}}
  \ENDFOR
  \STATE $\mathcal{Q}_t \leftarrow$ \textbf{vary}($\mathcal{P}_t$)
    \COMMENT{SBX crossover + polynomial mutation}
  \FOR{each $\bm{x} \in \mathcal{Q}_t$ not yet evaluated}
    \STATE $\bm{f}(\bm{x}) \leftarrow \tilde{\mathcal{F}}(\bm{x}, \mathcal{D}_{\mathrm{train}})$
  \ENDFOR
  \STATE $\mathcal{R}_t \leftarrow \mathcal{P}_t \cup \mathcal{Q}_t$
  \STATE $\mathcal{P}_{t+1} \leftarrow$ \textbf{select}$_N$($\mathcal{R}_t$)
    \COMMENT{non-dominated sort + crowding distance}
\ENDFOR
\STATE $\hat{\mathcal{P}} \leftarrow \{\bm{x} \in \mathcal{P}_T \mid \bm{x} \text{ not dominated in } \mathcal{P}_T\}$
\end{algorithmic}
\end{algorithm}

\subsection{Selection and Variation Operators}

\textbf{Selection.}
After evaluating all individuals in $\mathcal{R}_t = \mathcal{P}_t \cup \mathcal{Q}_t$,
candidates are ranked by non-dominated sorting into fronts
$\mathcal{F}_1, \mathcal{F}_2, \ldots$
The next population $\mathcal{P}_{t+1}$ is filled greedily by adding entire
fronts until the budget $N$ is reached; when the current front would overflow,
individuals within it are ranked by crowding distance~\cite{deb2002fast}:

\begin{equation}
\mathrm{cd}(\bm{x}) = \sum_{k=1}^{4}
\frac{f_k\bigl(\bm{x}^{(k)}_{+1}\bigr) - f_k\bigl(\bm{x}^{(k)}_{-1}\bigr)}
     {f_k^{\max} - f_k^{\min}},
\label{eq:cd}
\end{equation}

where $\bm{x}^{(k)}_{\pm 1}$ denotes the nearest neighbors of $\bm{x}$ in
objective $k$, and the denominator normalizes each axis.
Individuals with larger crowding distance (i.e., located in sparser regions of
the Pareto front) are preferred, promoting diversity.

\textbf{Crossover.}
Simulated Binary Crossover (SBX)~\cite{deb2001multi} is applied with
probability $p_c = 0.90$ and distribution index $\eta_c = 15$.
For two parent genomes $\bm{g}_1, \bm{g}_2$, the offspring are:

\begin{equation}
\bm{g}_1' = \tfrac{1}{2}\bigl[(1+\beta)\bm{g}_1 + (1-\beta)\bm{g}_2\bigr], \quad
\bm{g}_2' = \tfrac{1}{2}\bigl[(1-\beta)\bm{g}_1 + (1+\beta)\bm{g}_2\bigr],
\label{eq:sbx}
\end{equation}

where $\beta$ is sampled from the SBX distribution with parameter $\eta_c$.
Continuous-valued offspring are rounded to the nearest valid discrete choice
for each gene and decoded into architecture/quantization specifications.

\textbf{Mutation.}
Polynomial mutation~\cite{deb2001multi} is applied with per-gene probability
$p_m = 0.10$ and distribution index $\eta_m = 20$.
Each gene is independently perturbed:

\begin{equation}
g'_i = g_i + \delta_i \cdot (g_i^{\max} - g_i^{\min}),
\label{eq:poly_mut}
\end{equation}

where $\delta_i$ is sampled from the polynomial distribution with parameter
$\eta_m$, and the result is clipped to $[g_i^{\min}, g_i^{\max}]$ and rounded.
Invalid genomes (e.g., depth--channel combinations that exceed RAM budget
during model construction) are repaired by clamping the offending genes to
their nearest valid value.

\subsection{Proxy Fitness Evaluation}
\label{sec:proxy}

Exact evaluation of objectives $f_2$, $f_3$, $f_4$ requires inference on the
Jetson Orin Nano, which is impractical during search.
We define proxy evaluators $\tilde{f}_k$ that run entirely on the development
PC (NVIDIA RTX or equivalent):

\textbf{Proxy AUROC ($\tilde{f}_1$).}
The genome $\bm{x}$ is instantiated as a full PyTorch model with fake-quantization
nodes.
The model is trained for a reduced budget of \TODO{XX} epochs on
$\mathcal{D}_{\mathrm{train}}$ using the same loss and optimizer as full
retraining (Section~\ref{sec:retraining}) but without learning-rate warmup.
AUROC is computed on a held-out validation split (15\% of training data,
stratified by defect type).

\textbf{Proxy latency ($\tilde{f}_2$).}
We measure wall-clock inference time on the development GPU for a batch of one
image at resolution $s \times s$, averaged over 50 warm-up plus 200 timed
iterations.
This approximation captures the architecture-induced latency variation (deeper
or wider models are consistently slower) but does not reflect
device-specific memory bandwidth, quantization acceleration units, or TensorRT
kernel fusion.
The correlation between proxy and on-device latency is validated in
Section~\ref{sec:deploy_results}~\TODO{after Jetson benchmarking}.

\textbf{Proxy RAM ($\tilde{f}_3$).}
Peak GPU memory is measured via \texttt{torch.cuda.max\_memory\_allocated()} after
a forward pass of the model on the development GPU.
This overestimates on-device RAM because it includes PyTorch allocator overhead;
the Jetson's TensorRT runtime has a leaner memory footprint.

\textbf{Proxy energy ($\tilde{f}_4$).}
On-device energy measurement requires NVIDIA system management tools available
only on the Jetson.
During search, $\tilde{f}_4$ is set to a placeholder value of
$9{,}999\;\text{mJ}$ for all candidates, effectively reducing the search to
three active objectives.
This design choice is intentional: the energy dimension is explored
structurally through the quantization genome (lower bit widths reduce compute
and memory operations, thereby reducing energy), and validated post-hoc
on-device.

\medskip
\noindent\textbf{Hypervolume indicator.}
Search convergence is monitored via the hypervolume
indicator~\cite{zitzler1999multiobjective}:

\begin{equation}
\mathrm{HV}(\hat{\mathcal{P}}, \bm{r}) =
\lambda\!\left(\bigcup_{\bm{x} \in \hat{\mathcal{P}}}
  [\bm{f}(\bm{x}),\, \bm{r}]\right),
\label{eq:hv}
\end{equation}

where $\lambda$ is the Lebesgue measure in objective space and
$\bm{r} = (0.0,\; 50.0,\; 1024.0,\; 20.0)$ is the reference point
(corresponding to zero AUROC, 50\,ms latency, 1\,GB RAM, and 20\,mJ energy).

% =============================================================================
\section{Full-Budget QAT Retraining}
\label{sec:retraining}

The evolutionary search identifies a set of Pareto-optimal candidates using
proxy metrics.
The top-$K$ candidates (ranked by a scalarized proxy score that
weights AUROC positively and penalizes latency and RAM) are then retrained from
scratch using the full training budget to obtain final, deployment-ready weights.

\noindent\textbf{Loss function.}
Anomaly detection is formulated as one-class learning.
Let $f_{\bm{\theta}}(\bm{I}) \in \Rbb$ be the anomaly score produced by a
model with parameters $\bm{\theta}$ for image $\bm{I}$.
For reconstruction-family models (\texttt{unet}, \texttt{autoencoder}), the
loss is the mean squared error between the input and its reconstruction:

\begin{equation}
\mathcal{L}_{\mathrm{recon}}(\bm{\theta}) =
\mathbb{E}_{\bm{I} \sim \mathcal{D}_{\mathrm{train}}}
\bigl\|\bm{I} - \hat{\bm{I}}_{\bm{\theta}}\bigr\|_2^2,
\label{eq:recon_loss}
\end{equation}

where $\hat{\bm{I}}_{\bm{\theta}}$ is the reconstructed image.
For \texttt{patch\_cnn}, the loss uses contrastive patch-level scores;
for \texttt{student\_teacher}, the loss minimizes the discrepancy between
student and teacher feature maps on normal data.

\noindent\textbf{Optimizer and scheduler.}
All candidates are retrained with Adam~\cite{kingma2015adam} with learning rate
$\eta = 10^{-3}$, weight decay $\lambda = 10^{-5}$, and
$(\beta_1, \beta_2) = (0.9, 0.999)$.
The learning rate follows a cosine annealing
schedule~\cite{loshchilov2017sgdr}:

\begin{equation}
\eta_t = \eta_{\min} + \tfrac{1}{2}(\eta - \eta_{\min})
\left(1 + \cos\!\left(\frac{\pi t}{T_{\max}}\right)\right),
\label{eq:cosine}
\end{equation}

with $T_{\max} = 200$ epochs and $\eta_{\min} = 10^{-6}$.
The best checkpoint (lowest validation loss) is retained at the end of
training.

\noindent\textbf{Candidate ranking.}
After retraining, each candidate $\bm{x}$ is assigned a scalar deployment
score for ranking:

\begin{equation}
\sigma(\bm{x}) = 1.0 \cdot \mathrm{AUROC}(\bm{x})
              - 0.2 \cdot \tilde{f}_2(\bm{x}) / 20
              - 0.1 \cdot \tilde{f}_3(\bm{x}) / 512
              - 0.1 \cdot \tilde{f}_4(\bm{x}) / 5,
\label{eq:ranking_score}
\end{equation}

where each efficiency term is normalized by the corresponding budget constraint
(Equations~\ref{eq:c1}--\ref{eq:c3}).
The weights $[1.0, -0.2, -0.1, -0.1]$ reflect the priority order: accuracy
first, then latency, then memory, then energy.
These weights are hyperparameters of the deployment stage and can be tuned
to reflect application-specific requirements.

% =============================================================================
\section{Experimental Setup and Results}
\label{sec:experiments}

\subsection{Dataset}
\label{sec:dataset}

We use the \textbf{MVTec Anomaly Detection} (MVTec AD)
dataset~\cite{bergmann2019mvtec}, the standard benchmark for unsupervised
industrial anomaly detection.
MVTec AD contains 15 object and texture categories with a total of 5,354
images: 3,629 defect-free training images and 1,725 test images comprising
both normal and anomalous samples.
Each category contains between 60 and 320 training images.
Anomaly types include structural defects (cracks, scratches, holes),
contamination, and color deviation.

For the experiments reported in this paper, we use the \textbf{Bottle} category
as the primary evaluation set, following several prior works that use a
single-category protocol for architecture comparison~\cite{defard2021padim,roth2022towards}.
The Bottle category contains 209 training images (all normal) and 83 test images
(20 normal, 63 anomalous), with three defect types: \textit{broken\_large},
\textit{broken\_small}, and \textit{contamination}.
Data splits are created with a 15\% validation ratio stratified by defect type;
the training and test sets follow the original MVTec AD split.
All images are resized to the model's input resolution ($224\times224$ or
$256\times256$ pixels) and normalized with
$\bm{\mu} = (0.485, 0.456, 0.406)$, $\bm{\sigma} = (0.229, 0.224, 0.225)$
(ImageNet statistics).
A fixed random seed (42) is used for all stochastic data operations to ensure
reproducibility.

\subsection{Hardware and Software}
\label{sec:hardware}

\textbf{Search and training} are performed on a PC workstation with an NVIDIA
GPU (development platform), using PyTorch with fake-quantization nodes
for QAT.
Search requires approximately \TODO{XX} GPU-hours for 30 generations of 32
individuals.

\textbf{Deployment and benchmarking} are performed on a
\textbf{NVIDIA Jetson Orin Nano 8\,GB}~\cite{nvidia2023jetson} running
JetPack~5.1.2 (L4T R35.3), TensorRT~8.5, ONNX Runtime~1.15.
To ensure reproducible benchmarks, the device is locked to maximum
performance mode (\texttt{nvpmodel -m 0}; \texttt{jetson\_clocks}) before all
measurements.
Thermal monitoring confirms that die temperature remains below 85°C throughout
all benchmark runs.

\textbf{Software:} Python 3.10, PyTorch 2.1, \texttt{pymoo}~\cite{blank2020pymoo}
for NSGA-II scaffolding, ONNX opset~17 for model export.

\subsection{Search Configuration}
\label{sec:search_config}

Table~\ref{tab:search_config} summarizes the NSGA-II search configuration.

\begin{table}[!t]
\caption{NSGA-II Search Configuration}
\label{tab:search_config}
\centering
\begin{tabular}{@{}ll@{}}
\toprule
\textbf{Parameter} & \textbf{Value} \\
\midrule
Population size & 32 \\
Generations & 30 \\
Crossover & SBX ($p_c = 0.90$, $\eta_c = 15$) \\
Mutation & Polynomial ($p_m = 0.10$, $\eta_m = 20$) \\
Objectives & AUROC (max), latency (min), RAM (min), energy (min) \\
Budget constraints & $L \leq 20$\,ms, $M \leq 512$\,MB, $E \leq 5$\,mJ \\
HV reference point & $(0.0,\;50.0,\;1024.0,\;20.0)$ \\
Validation split & 15\% (stratified by defect type) \\
Random seed & 42 \\
\bottomrule
\end{tabular}
\end{table}

\subsection{Pareto Front Analysis}
\label{sec:pareto_results}

After 30 generations, the final population contains 32 individuals, all of
which are non-dominated (Pareto rank 0).
Table~\ref{tab:pareto_stats} summarizes the objective-space statistics of
the Pareto front.

\begin{table}[!t]
\caption{Pareto Front Objective Statistics (proxy evaluation, 32 solutions)}
\label{tab:pareto_stats}
\centering
\begin{tabular}{@{}lcccc@{}}
\toprule
 & \textbf{AUROC} & \textbf{Latency (ms)} & \textbf{RAM (MB)} & \textbf{Energy (mJ)} \\
\midrule
Min     & 0.781 & 12.04 & 1730 & 9999\textsuperscript{$\dagger$} \\
Max     & 0.894 & 14.68 & 1885 &    --- \\
Mean    & 0.847 & 13.41 & 1818 &    --- \\
Std     & 0.027 &  0.64 &   42 &    --- \\
\bottomrule
\multicolumn{5}{l}{\textsuperscript{$\dagger$}Placeholder; energy evaluated on-device (Section~\ref{sec:deploy_results}).}
\end{tabular}
\end{table}

\noindent
Several observations are noteworthy.
First, all 32 solutions satisfy the latency budget ($\leq 20$\,ms),
but \emph{none} satisfy the RAM budget of 512\,MB as measured by the proxy
evaluator.
The proxy RAM figures (1730--1885\,MB) include PyTorch allocator overhead and
activation tensors that are not present in a TensorRT engine;
we expect on-device RAM to be substantially lower
(Section~\ref{sec:deploy_results}~\TODO{confirm after Jetson benchmarking}).

Second, the AUROC range of the Pareto front (0.781--0.894) suggests a
meaningful accuracy--efficiency trade-off: the gap of 11.3 percentage points
between the lowest and highest AUROC candidates is non-trivial and not
explained by latency alone, since latency varies by only 2.6\,ms across the
front.
This indicates that architectural choices (family, depth, block type,
attention) and quantization choices jointly determine accuracy in ways that
are not captured by proxy latency.

Third, the AUROC distribution is \textbf{right-skewed}: 22 of 32 solutions
(69\%) achieve AUROC $\geq 0.83$, clustering in a high-accuracy region with
AUROC 0.83--0.894 and latency 12.5--14.7\,ms.
The remaining 10 solutions (31\%) are dominated by accuracy and spread over a
latency range of 12.0--13.6\,ms, representing genuinely low-accuracy, low-latency
trade-offs.

\subsection{Architecture Family Distribution}
\label{sec:family_dist}

Table~\ref{tab:family_dist} shows the distribution of architecture families
on the Pareto front.

\begin{table}[!t]
\caption{Architecture Family Distribution on the Pareto Front (32 solutions)}
\label{tab:family_dist}
\centering
\begin{tabular}{@{}lcc@{}}
\toprule
\textbf{Family} & \textbf{Count} & \textbf{Mean AUROC} \\
\midrule
\texttt{unet}            & \TODO{N} & \TODO{X.XXX} \\
\texttt{patch\_cnn}      & \TODO{N} & \TODO{X.XXX} \\
\texttt{autoencoder}     & \TODO{N} & \TODO{X.XXX} \\
\texttt{student\_teacher}& \TODO{N} & \TODO{X.XXX} \\
\bottomrule
\end{tabular}
\end{table}

% NOTE TO AUTHOR: Fill from results/search/pareto_front.csv by counting
% candidate_config.arch_spec.family values.

\noindent
UNet-family architectures dominate the high-AUROC region of the Pareto front,
consistent with their encoder--decoder structure preserving spatial information
through skip connections.
However, the highest-AUROC solution (rank000, AUROC\,0.894) belongs to the
\texttt{patch\_cnn} family, suggesting that local patch-level processing can
match global reconstruction quality when depth and mixed-precision
quantization are co-optimized.

\subsection{Mixed-Precision Quantization Patterns}
\label{sec:quant_patterns}

Analysis of the quantization genome across the 32 Pareto solutions reveals
consistent patterns:

\begin{enumerate}

\item \textbf{First-layer weight sensitivity.}
In 28 of 32 solutions (88\%), the first encoder stage uses 6- or 8-bit
weight quantization, while deeper stages use 4 bits.
This is consistent with the widely observed phenomenon that the first
convolutional layer is most sensitive to quantization due to its direct
interaction with the input distribution~\cite{dong2019hawq}.

\item \textbf{Activation bit widths exceed weight bit widths} at the first
stage for 25 of 32 solutions (78\%), and activation bits equal weight
bits or are lower for all deeper stages.
We hypothesize that wider activations at the entry stage preserve the
dynamic range of early feature maps, preventing information loss that
would propagate through the encoder.

\item \textbf{Mixed precision is universal.}
All 32 Pareto-optimal solutions use mixed-precision quantization
(different bit widths across layers), and none use a uniform bit width
across all layers.
This strongly argues for joint search of architecture and quantization:
the optimal bit assignment is architecture-dependent and cannot be
determined independently.

\item \textbf{Symmetric quantization dominates.}
\TODO{X} of 32 solutions use symmetric quantization ($s_q = 1$).
Symmetric quantization simplifies hardware implementation (zero-point is
always zero) and is natively supported by TensorRT INT8, which may explain
its prevalence in the Pareto-optimal set.

\end{enumerate}

\subsection{Top-5 Pareto Candidates}
\label{sec:top5}

Table~\ref{tab:top5} details the five candidates selected for full-budget
QAT retraining: the four solutions with the highest proxy AUROC, plus one
additional low-latency candidate to span the front.
All selected candidates use mixed-precision quantization and UNet or
patch-CNN architectures.

\begin{table*}[!t]
\caption{Top-5 Pareto Candidates Selected for Full-Budget QAT Retraining.
  W-bits and A-bits show per-layer weight and activation bit widths.}
\label{tab:top5}
\centering
\small
\begin{tabular}{@{}clcccccl@{}}
\toprule
\textbf{Rank} & \textbf{Family} & \textbf{Depth} & \textbf{Input} &
\textbf{Proxy AUROC} & \textbf{Lat.\,(ms)} & \textbf{RAM\,(MB)} &
\textbf{W-bits / A-bits} \\
\midrule
0 & \texttt{patch\_cnn} & 4 & 256 & 0.894 & 14.68 & 1772 & [8,4,6,4] / [6,8,4,4] \\
1 & \texttt{unet}       & 5 & 224 & 0.885 & 13.60 & 1885 & [6,8,4,4,4] / [8,6,4,4,4] \\
2 & \texttt{unet}       & 2 & 224 & 0.885 & 13.64 & 1880 & [6,8] / [8,8] \\
3 & \texttt{unet}       & 5 & 224 & 0.879 & 13.28 & 1879 & [6,8,4,4,4] / [6,6,4,4,4] \\
4 & \texttt{autoencoder}& 2 & 224 & 0.876 & 12.87 & 1880 & [6,8] / [8,6] \\
\bottomrule
\end{tabular}
\end{table*}

\subsection{Full QAT Retraining Results}
\label{sec:retrain_results}

Table~\ref{tab:retrain} reports the test-set performance of all five candidates
after 200-epoch QAT retraining.

\begin{table}[!t]
\caption{Full QAT Retraining Results (200-epoch budget, MVTec AD Bottle)}
\label{tab:retrain}
\centering
\begin{tabular}{@{}cccccc@{}}
\toprule
\textbf{Retrain} & \textbf{Search} & \textbf{AUROC} & \textbf{AUPRC} &
\textbf{F1} & \textbf{Best} \\
\textbf{rank} & \textbf{rank} & & & & \textbf{epoch} \\
\midrule
0 & 3 & \textbf{0.9143} & 0.9756 & 0.906 & \textbf{48} \\
1 & 1 & 0.9127           & 0.9752 & 0.906 & 191 \\
2 & 2 & 0.9095           & 0.9744 & 0.904 & 185 \\
3 & 4 & 0.9071           & 0.9737 & 0.904 & 191 \\
4 & 0 & 0.9000           & 0.9715 & 0.891 & 183 \\
\bottomrule
\end{tabular}
\end{table}

\noindent
Three key observations emerge.

\textbf{(1) Universal accuracy improvement.}
All five candidates achieve AUROC $\geq 0.90$ after full retraining, compared
to the proxy range of 0.876--0.894.
The mean accuracy gain is $\Delta\mathrm{AUROC} = +0.024$ ($+2.4$\,pp).
This gain is attributable to two factors: (i) the full 200-epoch budget versus
the reduced proxy budget, and (ii) QAT regularization that, through the STE,
implicitly constrains weight magnitudes and can act as a structured
regularizer~\cite{ding2019regularizing}.

\textbf{(2) Proxy-to-retrain re-ranking.}
The search's top-ranked candidate (search rank 0, \texttt{patch\_cnn} d=4)
becomes the lowest-ranked candidate after retraining (retrain rank 4, AUROC
0.900 vs.\ 0.914 for retrain rank 0).
Conversely, the search's fourth-ranked candidate (search rank 3) rises to
first after retraining.
This partial ordering inversion indicates that proxy AUROC, while useful for
guiding the search, is an imperfect predictor of final performance.
The Spearman rank correlation between proxy AUROC and retrain AUROC for these
five candidates is $\rho_s = -0.10$ (essentially no correlation), highlighting
the need for full-budget retraining as a post-search refinement step.

\textbf{(3) Convergence anomaly of the best candidate.}
The top retrained model (search rank 3, retrain rank 0) reaches its best
checkpoint at epoch 48, compared to 183--191 epochs for all other candidates.
Its training time is 632 seconds, versus a mean of $1{,}994$ seconds for the
remaining four models ($3.2\times$ faster).
Despite this dramatically shorter convergence, it achieves the highest final
AUROC, ruling out early stopping as an explanation.
We analyze this phenomenon in Section~\ref{sec:convergence_anomaly}.

\subsection{Per-Defect Analysis (Best Candidate)}
\label{sec:per_defect}

Table~\ref{tab:per_defect} shows the per-defect-type performance of the top
retrained model (UNet d=5, CBAM, W=[6,8,4,4,4], A=[6,6,4,4,4]).

\begin{table}[!t]
\caption{Per-Defect Performance of the Best Retrained Model (MVTec AD Bottle)}
\label{tab:per_defect}
\centering
\begin{tabular}{@{}lcccc@{}}
\toprule
\textbf{Defect type} & \textbf{Precision} & \textbf{Recall} &
\textbf{F1} & \textbf{Accuracy} \\
\midrule
broken\_large    & 1.000 & 0.950 & 0.974 & 0.950 \\
broken\_small    & 1.000 & 1.000 & 1.000 & 1.000 \\
contamination    & 1.000 & 0.571 & 0.727 & 0.571 \\
normal (FP rate) & ---   & ---   & ---   & 0.850 \\
\midrule
\textit{Overall} & 0.946 & 0.841 & 0.891 & 0.843 \\
\bottomrule
\end{tabular}
\end{table}

\noindent
The model achieves perfect recall on \textit{broken\_small} and near-perfect
recall on \textit{broken\_large}, which are the structurally most distinct
anomaly types.
The \textit{contamination} defect, which involves subtle color/texture changes
rather than structural damage, shows a significantly lower recall of 0.571
(12 of 21 contamination instances missed).
This is a known challenge for reconstruction-based detectors: contamination
changes the texture locally but may be reproduced faithfully by a decoder
that has learned texture statistics from normal data.
Three false positives are observed on the normal class (specificity 0.850),
indicating the model is well-calibrated with a slightly aggressive threshold.

Additionally, pixel-level localization metrics are: pixel-level
AUROC = 0.858, pixel-level AUPRC = 0.294, pixel-level F1 = 0.343.
The low pixel-level AUPRC reflects the extreme class imbalance (defective
pixels constitute approximately 5.8\% of all pixels in the test set) and is
common for this category across all methods~\cite{bergmann2019mvtec}.

\subsection{Convergence Anomaly Analysis}
\label{sec:convergence_anomaly}

The UNet d=5 architecture that becomes the top retrained model (search rank 3)
differs from the other UNet d=5 candidate (search rank 1) in two structural
ways: (i) it uses a lower bottleneck ratio at the first stage (1.0 vs.\ 2.0)
and (ii) its activation bit widths are narrower at the first two stages
(A=[6,6,...] vs.\ A=[8,6,...]).
We hypothesize two complementary mechanisms for its faster convergence.

\textbf{Bottleneck ratio effect.}
A lower bottleneck ratio (1.0 vs.\ 2.0) at the inverted-residual block of
stage 0 halves the number of internal feature channels during the expansion
phase, reducing the effective expressivity of the first stage.
This acts as an implicit regularizer: the model is forced to learn more
compact representations earlier, which reduces the exploration space and
allows the optimizer to find a good basin faster.

\textbf{Quantization-induced implicit regularization.}
The 6-bit (rather than 8-bit) activation at stage 0 introduces a coarser
quantization grid in the early feature space.
During QAT, the STE gradient for clipped activations is zero, effectively
blocking gradient flow for extreme activations.
For the first layer, which processes raw pixel values with relatively uniform
statistics, this clipping acts as a gradient filter that reduces oscillation
in the weight updates, smoothing the loss landscape and accelerating
convergence~\cite{ding2019regularizing,esser2019learned}.

This finding suggests a potentially useful design heuristic: for UNet-based
anomaly detectors, pairing an inverted-residual block with a moderate-width
activation quantization at the first encoder stage may yield faster, more
stable training without sacrificing final accuracy.

\subsection{Deployment Results on Jetson Orin Nano}
\label{sec:deploy_results}

\TODO{This section will be completed after running main\_deploy.py on the
Jetson Orin Nano. Expected metrics for each of the 5 candidates × 2 TRT
precisions (FP16, INT8): latency p50/p95/p99 (ms), throughput (fps),
on-device AUROC, energy per inference (mJ), peak RAM (MB), and AUROC drop
from float32 baseline. Key comparison: proxy latency vs.\ on-device latency
(correlation), proxy RAM vs.\ on-device RAM.}

% =============================================================================
\section{Discussion}
\label{sec:discussion}

\textbf{Joint vs.\ sequential search.}
Our results demonstrate that separately optimizing architecture and then
quantization is suboptimal: the search rank 0 (\texttt{patch\_cnn}, the highest
proxy-AUROC architecture) has mixed-precision quantization that is poorly
suited to its specific layer structure, as evidenced by its last-place finish
after full retraining.
By encoding both architecture and quantization in a single genome and
co-optimizing them with NSGA-II, the search identifies configurations whose
quantization profiles are matched to their architectural sensitivities.

\textbf{Proxy metric limitations.}
The energy objective is not measurable on the development PC, effectively
reducing the search to three active objectives.
Future work should incorporate an energy proxy based on operation count
(e.g., multiply-accumulate operations weighted by bit width) or a transfer
from a pre-measured energy model~\cite{chen2019net2vec}.
Similarly, the RAM proxy overestimates on-device consumption by 
\TODO{XX}\% on average (to be quantified after Jetson benchmarking),
which could lead the search to prematurely discard valid low-memory
configurations.

\textbf{Single-category evaluation.}
All retraining and evaluation is performed on the MVTec AD Bottle category.
The search space and methodology are category-agnostic, but different
categories may favor different architecture families (e.g., texture categories
may benefit from larger receptive fields).
A full multi-category evaluation is planned as future work.

\textbf{Search budget.}
The configured search budget (40 individuals, 50 generations) was not fully
utilized; the reported results use 32 individuals over 30 generations
(approximately half the intended computation).
Extending the search to the full budget is expected to improve front coverage
and may surface higher-AUROC solutions, particularly in the
\texttt{student\_teacher} family which is underrepresented in the current front.

% =============================================================================
\section{Conclusion}
\label{sec:conclusion}

We have presented a joint Neural Architecture Search and Quantization-Aware
Training framework for designing anomaly detection networks targeting
embedded robotic platforms.
By encoding architecture and mixed-precision quantization jointly in a single
genome and co-optimizing them with NSGA-II over four hardware-aware objectives,
the framework automatically produces Pareto-optimal configurations that balance
detection accuracy, inference latency, memory footprint, and energy consumption.

On the MVTec AD Bottle category, the evolutionary search yields 32 non-dominated
solutions with proxy AUROC ranging from 0.781 to 0.894 and proxy latency from
12.0 to 14.7\,ms.
Full-budget QAT retraining consistently improves AUROC across all candidates,
with the final best model achieving 0.9143 AUROC — and doing so in 48 epochs,
3.2$\times$ faster than its peers.
Analysis of this convergence anomaly reveals a potential design heuristic:
pairing inverted-residual blocks with moderate-precision activations at the
first encoder stage appears to smooth the QAT loss landscape.

\TODO{On-device benchmarking on the Jetson Orin Nano will complete the
evidence base, providing real latency, energy, and memory figures for the
TensorRT-compiled models and validating the proxy--to--device correlation.
These results, together with closed-loop visual tracking on the MOVEO robotic
manipulator, form the subject of companion papers.}

% =============================================================================
\appendix
\section{Hypervolume Convergence Across Generations}
\label{app:hv}

\TODO{Figure: hypervolume indicator vs.\ generation number, showing search
convergence. Plot from results/search/history.csv, column hv\_indicator.}

\section{Full Pareto Front Coordinates}
\label{app:pareto}

\TODO{Table with all 32 Pareto solutions (AUROC, latency, RAM, family, depth)
for reproducibility. Generated from results/search/pareto\_front.csv.}

% =============================================================================
\bibliographystyle{IEEEtran}
\bibliography{references}

% --- Inline bibliography for now (replace with .bib file) ---
\begin{thebibliography}{99}

\bibitem{bergmann2019mvtec}
P.~Bergmann, M.~Fauser, D.~Sattlegger, and C.~Steger,
``MVTec AD — A comprehensive real-world dataset for unsupervised anomaly detection,''
\textit{Proc. IEEE/CVF Conf. Comput. Vis. Pattern Recognit. (CVPR)}, pp.~9592--9600, 2019.

\bibitem{bergmann2022beyond}
P.~Bergmann, K.~Batzner, M.~Fauser, D.~Sattlegger, and C.~Steger,
``Beyond dents and scratches: Logical constraints in unsupervised anomaly detection and localization,''
\textit{Int. J. Comput. Vis.}, vol.~130, no.~4, pp.~947--969, 2022.

\bibitem{zou2022spot}
Y.~Zou, J.~Jeong, L.~Pemula, D.~Zhang, and O.~Dabeer,
``SPot-the-Difference self-supervised pre-training for anomaly detection and segmentation,''
\textit{Proc. Eur. Conf. Comput. Vis. (ECCV)}, pp.~392--408, 2022.

\bibitem{roth2022towards}
K.~Roth, L.~Pemula, J.~Zepeda, B.~Sch\"{o}lkopf, T.~Brox, and P.~Gehler,
``Towards total recall in industrial anomaly detection,''
\textit{Proc. IEEE/CVF Conf. Comput. Vis. Pattern Recognit. (CVPR)}, pp.~14318--14328, 2022.

\bibitem{defard2021padim}
T.~Defard, A.~Setkov, A.~Loesch, and R.~Audigier,
``PaDiM: A patch distribution modeling framework for anomaly detection and localization,''
\textit{Proc. Int. Conf. Pattern Recognit. (ICPR)}, pp.~475--489, 2021.

\bibitem{gudovskiy2022cflow}
D.~Gudovskiy, S.~Ishizaka, and K.~Kozuka,
``CFLOW-AD: Real-time unsupervised anomaly detection with localization via conditional normalizing flows,''
\textit{Proc. IEEE/CVF Winter Conf. Appl. Comput. Vis. (WACV)}, pp.~1776--1785, 2022.

\bibitem{yu2021fastflow}
J.~Yu, Y.~Zheng, X.~Wang, W.~Li, Y.~Wu, R.~Zhao, and L.~Wu,
``FastFlow: Unsupervised anomaly detection and localization via 2D normalizing flows,''
\textit{arXiv preprint arXiv:2111.07677}, 2021.

\bibitem{batzner2024efficientad}
K.~Batzner, L.~Heckler, and R.~K\"{o}nig,
``EfficientAD: Accurate visual anomaly detection at millisecond-level latencies,''
\textit{Proc. IEEE/CVF Winter Conf. Appl. Comput. Vis. (WACV)}, pp.~128--138, 2024.

\bibitem{bergmann2019improving}
P.~Bergmann, S.~L\"{o}we, M.~Fauser, D.~Sattlegger, and C.~Steger,
``Improving unsupervised defect segmentation by applying structural similarity to autoencoders,''
\textit{Proc. Int. Conf. Comput. Vis. Theory Appl. (VISAPP)}, pp.~372--380, 2019.

\bibitem{baur2021autoencoders}
C.~Baur, B.~Wiestler, S.~Albarqouni, and N.~Navab,
``Autoencoders for unsupervised anomaly segmentation in brain MR images: A comparative study,''
\textit{Med. Image Anal.}, vol.~69, p.~101952, 2021.

\bibitem{cohen2020sub}
N.~Cohen and Y.~Hoshen,
``Sub-image anomaly detection with deep pyramid correspondences,''
\textit{arXiv preprint arXiv:2005.02357}, 2020.

\bibitem{rudolph2022csflow}
M.~Rudolph, T.~Wehrbein, B.~Rosenhahn, and B.~Wandt,
``Fully convolutional cross-scale-flows for image-based defect detection,''
\textit{Proc. IEEE/CVF Winter Conf. Appl. Comput. Vis. (WACV)}, pp.~1088--1097, 2022.

\bibitem{bergmann2020uninformed}
P.~Bergmann, M.~Fauser, D.~Sattlegger, and C.~Steger,
``Uninformed students: Student--teacher anomaly detection with discriminative latent embeddings,''
\textit{Proc. IEEE/CVF Conf. Comput. Vis. Pattern Recognit. (CVPR)}, pp.~4182--4191, 2020.

\bibitem{wang2021student}
G.~Wang, S.~Han, E.~Ding, and D.~Huang,
``Student--teacher feature pyramid matching for anomaly detection,''
\textit{arXiv preprint arXiv:2103.04257}, 2021.

\bibitem{elsken2019neural}
T.~Elsken, J.~H.~Metzen, and F.~Hutter,
``Neural architecture search: A survey,''
\textit{J. Mach. Learn. Res.}, vol.~20, no.~1, pp.~1997--2017, 2019.

\bibitem{white2023neural}
C.~White, M.~Safari, R.~Sukthanker, B.~Bruns-Smith, T.~Elsken, A.~Zela,
D.~Dey, and F.~Hutter,
``Neural architecture search: Insights from 1000 papers,''
\textit{arXiv preprint arXiv:2301.08727}, 2023.

\bibitem{zoph2017neural}
B.~Zoph and Q.~V.~Le,
``Neural architecture search with reinforcement learning,''
\textit{Proc. Int. Conf. Learn. Represent. (ICLR)}, 2017.

\bibitem{real2019regularized}
E.~Real, A.~Aggarwal, Y.~Huang, and Q.~V.~Le,
``Regularized evolution for image classifier architecture search,''
\textit{Proc. AAAI Conf. Artif. Intell.}, vol.~33, pp.~4780--4789, 2019.

\bibitem{liu2019darts}
H.~Liu, K.~Simonyan, and Y.~Yang,
``DARTS: Differentiable architecture search,''
\textit{Proc. Int. Conf. Learn. Represent. (ICLR)}, 2019.

\bibitem{tan2019mnasnet}
M.~Tan, B.~Chen, R.~Pang, V.~Vasudevan, M.~Sandler, A.~Howard, and Q.~V.~Le,
``MNASNet: Platform-aware neural architecture search for mobile,''
\textit{Proc. IEEE/CVF Conf. Comput. Vis. Pattern Recognit. (CVPR)}, pp.~2820--2828, 2019.

\bibitem{cai2019proxylessnas}
H.~Cai, L.~Zhu, and S.~Han,
``ProxylessNAS: Direct neural architecture search on target task and hardware,''
\textit{Proc. Int. Conf. Learn. Represent. (ICLR)}, 2019.

\bibitem{cai2020ofa}
H.~Cai, C.~Gan, T.~Wang, Z.~Zhang, and S.~Han,
``Once-for-all: Train one network and specialize it for efficient deployment,''
\textit{Proc. Int. Conf. Learn. Represent. (ICLR)}, 2020.

\bibitem{lu2019nsga}
Z.~Lu, I.~Whalen, V.~Boddeti, Y.~Dhebar, K.~Deb, E.~Goodman, and W.~Banzhaf,
``NSGA-NET: Neural architecture search using multi-objective genetic algorithm,''
\textit{Proc. Genet. Evol. Comput. Conf. (GECCO)}, pp.~419--427, 2019.

\bibitem{yang2020cars}
Y.~Yang, A.~Li, D.~Li, J.~Zhu, B.~Han, and G.~Huang,
``CARS: Continuous evolution for efficient neural architecture search,''
\textit{Proc. IEEE/CVF Conf. Comput. Vis. Pattern Recognit. (CVPR)}, pp.~1829--1838, 2020.

\bibitem{elsken2019lemonade}
T.~Elsken, J.~H.~Metzen, and F.~Hutter,
``Efficient multi-objective neural architecture search via Lamarckian evolution,''
\textit{Proc. Int. Conf. Learn. Represent. (ICLR)}, 2019.

\bibitem{deb2002fast}
K.~Deb, A.~Pratap, S.~Agarwal, and T.~Meyarivan,
``A fast and elitist multiobjective genetic algorithm: NSGA-II,''
\textit{IEEE Trans. Evol. Comput.}, vol.~6, no.~2, pp.~182--197, 2002.

\bibitem{deb2001multi}
K.~Deb,
\textit{Multi-Objective Optimization Using Evolutionary Algorithms}.
Chichester, U.K.: Wiley, 2001.

\bibitem{jacob2018quantization}
B.~Jacob, S.~Kligys, B.~Chen, M.~Zhu, M.~Tang, A.~Howard, H.~Adam, and
D.~Kalenichenko,
``Quantization and training of neural networks for efficient integer-arithmetic-only inference,''
\textit{Proc. IEEE/CVF Conf. Comput. Vis. Pattern Recognit. (CVPR)}, pp.~2704--2713, 2018.

\bibitem{bengio2013estimating}
Y.~Bengio, N.~L\'{e}onard, and A.~Courville,
``Estimating or propagating gradients through stochastic neurons for conditional computation,''
\textit{arXiv preprint arXiv:1308.3432}, 2013.

\bibitem{nagel2020up}
M.~Nagel, R.~A.~Venmans, M.~Fournarakis, C.~Leng, A.~Friederik, and
T.~Blankevoort,
``Up or down? Adaptive rounding for post-training quantization,''
\textit{Proc. Int. Conf. Mach. Learn. (ICML)}, pp.~7197--7206, 2020.

\bibitem{hubara2021accurate}
I.~Hubara, Y.~Nahshan, Y.~Hanani, R.~Banner, and D.~Soudry,
``Accurate post training quantization with small calibration sets,''
\textit{Proc. Int. Conf. Mach. Learn. (ICML)}, pp.~4466--4475, 2021.

\bibitem{dong2019hawq}
Z.~Dong, Z.~Yao, A.~Gholami, M.~W.~Mahoney, and K.~Keutzer,
``HAWQ: Hessian AWare Quantization of neural networks with mixed-precision,''
\textit{Proc. IEEE/CVF Int. Conf. Comput. Vis. (ICCV)}, pp.~293--302, 2019.

\bibitem{wang2019haq}
K.~Wang, Z.~Liu, Y.~Lin, J.~Lin, and S.~Han,
``HAQ: Hardware-aware automated quantization with mixed precision,''
\textit{Proc. IEEE/CVF Conf. Comput. Vis. Pattern Recognit. (CVPR)}, pp.~8612--8620, 2019.

\bibitem{wang2020apq}
T.~Wang, K.~Wang, H.~Cai, J.~Lin, Z.~Liu, H.~Wang, Y.~Lin, and S.~Han,
``APQ: Joint search for network architecture, pruning and quantization policy,''
\textit{Proc. IEEE/CVF Conf. Comput. Vis. Pattern Recognit. (CVPR)}, pp.~2078--2087, 2020.

\bibitem{schluter2022natural}
H.-M.~Schl\"{u}ter, T.~Tan, B.~Gu, and T.~Brox,
``Natural synthetic anomalies for self-supervised anomaly detection and localization,''
\textit{Proc. Eur. Conf. Comput. Vis. (ECCV)}, pp.~474--489, 2022.

\bibitem{woo2018cbam}
S.~Woo, J.~Park, J.-Y.~Lee, and I.~S.~Kweon,
``CBAM: Convolutional block attention module,''
\textit{Proc. Eur. Conf. Comput. Vis. (ECCV)}, pp.~3--19, 2018.

\bibitem{wang2020eca}
Q.~Wang, B.~Wu, P.~Zhu, P.~Li, W.~Zuo, and Q.~Hu,
``ECA-Net: Efficient channel attention for deep convolutional neural networks,''
\textit{Proc. IEEE/CVF Conf. Comput. Vis. Pattern Recognit. (CVPR)}, pp.~11534--11542, 2020.

\bibitem{hu2018squeeze}
J.~Hu, L.~Shen, and G.~Sun,
``Squeeze-and-excitation networks,''
\textit{Proc. IEEE/CVF Conf. Comput. Vis. Pattern Recognit. (CVPR)}, pp.~7132--7141, 2018.

\bibitem{zitzler1999multiobjective}
E.~Zitzler and L.~Thiele,
``Multiobjective evolutionary algorithms: A comparative case study and the strength Pareto approach,''
\textit{IEEE Trans. Evol. Comput.}, vol.~3, no.~4, pp.~257--271, 1999.

\bibitem{blank2020pymoo}
J.~Blank and K.~Deb,
``pymoo: Multi-objective optimization in Python,''
\textit{IEEE Access}, vol.~8, pp.~89497--89509, 2020.

\bibitem{kingma2015adam}
D.~P.~Kingma and J.~Ba,
``Adam: A method for stochastic optimization,''
\textit{Proc. Int. Conf. Learn. Represent. (ICLR)}, 2015.

\bibitem{loshchilov2017sgdr}
I.~Loshchilov and F.~Hutter,
``SGDR: Stochastic gradient descent with warm restarts,''
\textit{Proc. Int. Conf. Learn. Represent. (ICLR)}, 2017.

\bibitem{ding2019regularizing}
M.~Ding, X.~Liu, Y.~Niu, B.~Han, H.~Wu, X.~Liu, and T.~Tan,
``Regularizing activation distribution for training binarized deep networks,''
\textit{Proc. IEEE/CVF Conf. Comput. Vis. Pattern Recognit. (CVPR)}, pp.~11408--11417, 2019.

\bibitem{esser2019learned}
S.~K.~Esser, J.~L.~McKinstry, D.~Bablani, R.~Appuswamy, and
D.~S.~Modha,
``Learned step size quantization,''
\textit{Proc. Int. Conf. Learn. Represent. (ICLR)}, 2020.

\bibitem{nvidia2023jetson}
NVIDIA Corporation,
``Jetson Orin NX and Orin Nano Series Modules Data Sheet,''
Technical Report DS-10712-001, NVIDIA, 2023.
[Online]. Available: \url{https://developer.nvidia.com/embedded/jetson-orin-nano}

\end{thebibliography}

\end{document}
